\chapter{Geometry and Dynamics of Learning}

\section{Parameter Space and Function Space}

Training a neural network can be viewed from two closely related perspectives:
\begin{itemize}
    \item \textbf{parameter space:} the space of weights and biases,
    \item \textbf{function space:} the space of input--output maps realized by the network.
\end{itemize}

\medskip

\noindent
If the network is parameterized by \(\theta \in \mathbb{R}^d\), then training minimizes a loss
\[
L(\theta) = \frac{1}{n}\sum_{i=1}^n \ell\bigl(f_\theta(x_i),y_i\bigr),
\]
where \(f_\theta\) is the function represented by the network.

\medskip

\noindent
\textbf{Parameter-space view.}
\begin{itemize}
    \item A point \(\theta\) represents one specific choice of parameters
    \item Optimization moves through a high-dimensional landscape \(L(\theta)\)
\end{itemize}

\medskip

\noindent
\textbf{Function-space view.}
\begin{itemize}
    \item The object of real interest is not the parameter vector itself, but the function \(f_\theta\)
    \item Different parameter vectors may represent similar or even identical functions
\end{itemize}

\medskip

\noindent
This distinction is important because neural networks are often highly redundant:
\begin{itemize}
    \item permutations of hidden units can leave the realized function unchanged,
    \item rescalings in some architectures can produce equivalent predictions,
    \item large motions in parameter space can correspond to small motions in function space.
\end{itemize}

\medskip

\noindent
\textbf{Insight.}
Optimization takes place in parameter space, but generalization is about the learned function. Much of modern theory studies how parameter dynamics induce trajectories in function space.

\section{Loss Landscapes}

The loss function defines a high-dimensional landscape over parameter space:
\[
L(\theta): \mathbb{R}^d \to \mathbb{R}.
\]

\medskip

\noindent
\textbf{Landscape interpretation.}
\begin{itemize}
    \item each point corresponds to a parameter configuration,
    \item the height of the landscape is the loss value,
    \item learning corresponds to moving downhill.
\end{itemize}

\medskip

\noindent
Although this picture is easy to draw in low dimensions, the true loss landscape of a neural network is extremely high-dimensional and far more intricate.

\medskip

\noindent
\textbf{Critical points.}
A point \(\theta^\ast\) is critical if
\[
\nabla L(\theta^\ast)=0.
\]
Among critical points, one distinguishes:
\begin{itemize}
    \item \textbf{local minima:} directions of nonnegative curvature,
    \item \textbf{local maxima:} directions of nonpositive curvature,
    \item \textbf{saddle points:} some directions curve upward and others downward.
\end{itemize}

\medskip

\noindent
In high dimensions, saddle points are often more prevalent than poor isolated local minima. This helps explain why optimization difficulties in deep learning are not well described simply as ``getting stuck in bad local minima.''

\medskip

\noindent
\textbf{Flat and sharp regions.}
\begin{itemize}
    \item \textbf{Flat regions} have small curvature in many directions
    \item \textbf{Sharp regions} have large curvature, so loss changes rapidly under perturbations
\end{itemize}

\medskip

\noindent
This matters because optimization speed, stability, and sometimes generalization behavior are strongly influenced by local geometry.

\medskip

\noindent
\textbf{Caution.}
In neural networks, geometric notions such as flatness can depend on parametrization, so they must be interpreted carefully. Still, the landscape picture remains a useful organizing idea.

\section{Curvature, Hessians, and Optimization Geometry}

To move beyond first-order descent intuition, one studies curvature.

\subsection{Hessian Matrix}

The local second-order geometry of the loss is captured by the Hessian:
\[
H(\theta)=\nabla^2 L(\theta).
\]

\medskip

\noindent
\textbf{Interpretation of the Hessian.}
\begin{itemize}
    \item it measures how the gradient changes locally,
    \item it describes curvature in different directions,
    \item it helps distinguish minima, saddles, and flat directions.
\end{itemize}

\medskip

\noindent
For a direction \(v \in \mathbb{R}^d\), the quadratic approximation gives
\[
L(\theta + v) \approx L(\theta) + \nabla L(\theta)^\top v + \frac12 v^\top H(\theta) v.
\]
Thus the quantity
\[
v^\top H(\theta) v
\]
describes the curvature of the loss in direction \(v\).

\subsection{Eigenvalue Intuition}

The eigenvalues of the Hessian give local curvature scales.

\begin{itemize}
    \item Large positive eigenvalues indicate steep upward curvature
    \item Small eigenvalues indicate nearly flat directions
    \item Negative eigenvalues indicate descent directions and hence saddle-like structure
\end{itemize}

\medskip

\noindent
\textbf{Hessian intuition.}
Neural-network losses often exhibit:
\begin{itemize}
    \item a few directions of strong curvature,
    \item many weakly curved or nearly flat directions,
    \item occasional negative-curvature directions away from minima.
\end{itemize}

\medskip

\noindent
This suggests that optimization takes place in a highly anisotropic geometry rather than an isotropic bowl.

\subsection{Curvature and Optimization}

Curvature matters because gradient descent behaves differently in different directions.

\medskip

\noindent
If one applies the update
\[
\theta_{k+1}=\theta_k-\eta \nabla L(\theta_k),
\]
then large curvature directions constrain the step size \(\eta\). If \(\eta\) is too large relative to local curvature, the iteration can oscillate or diverge.

\medskip

\noindent
\textbf{Consequences of curvature.}
\begin{itemize}
    \item steep directions slow stable optimization,
    \item flat directions lead to slow progress,
    \item anisotropy motivates momentum, preconditioning, and adaptive methods.
\end{itemize}

\medskip

\noindent
Second-order methods explicitly try to exploit Hessian information, but in large neural networks exact Hessian computations are often too expensive. Still, Hessian-based intuition is central even when one uses only first-order optimizers.

\subsection{Geometry Beyond the Hessian}

The Hessian is a local object. Deep learning also involves global geometric features:
\begin{itemize}
    \item connected low-loss valleys,
    \item broad basins,
    \item ridges and plateaus,
    \item symmetries and degeneracies of parametrization.
\end{itemize}

\medskip

\noindent
Thus optimization geometry is not merely about one critical point, but about the structure of large low-loss regions and the trajectories connecting them.

\section{Gradient Flow and Dynamical Systems}

Optimization can be studied not only as an iterative algorithm but also as a dynamical system.

\subsection{From Gradient Descent to Gradient Flow}

Discrete-time gradient descent takes the form
\[
\theta_{k+1}=\theta_k-\eta \nabla L(\theta_k).
\]
When the step size \(\eta\) is small, one is led to the continuous-time limit
\[
\frac{d\theta(t)}{dt} = -\nabla L(\theta(t)),
\]
called \textbf{gradient flow}.

\medskip

\noindent
\textbf{Interpretation.}
The parameter vector evolves in the direction of steepest instantaneous decrease of the loss.

\subsection{Energy Dissipation}

Along a gradient-flow trajectory, the loss decreases monotonically:
\[
\frac{d}{dt}L(\theta(t))
=
\nabla L(\theta(t))^\top \frac{d\theta(t)}{dt}
=
-\|\nabla L(\theta(t))\|^2
\le 0.
\]

\medskip

\noindent
This identity is fundamental. It shows that gradient flow is a dissipative dynamical system: the loss acts like an energy that decreases over time.

\subsection{Fixed Points and Stability}

A fixed point of gradient flow satisfies
\[
\frac{d\theta(t)}{dt}=0,
\]
so necessarily
\[
\nabla L(\theta)=0.
\]
Thus critical points of the loss correspond to equilibria of the flow.

\medskip

\noindent
Their stability depends on curvature:
\begin{itemize}
    \item minima are locally attracting under suitable conditions,
    \item saddle points are unstable in directions of negative curvature,
    \item flat directions can produce slow or non-isolated behavior.
\end{itemize}

\medskip

\noindent
This dynamical-systems viewpoint explains why negative curvature helps optimization escape many saddles.

\subsection{Stochastic Dynamics}

Stochastic gradient descent introduces noise because gradients are estimated from minibatches. At a high level, one models this by a stochastic differential equation
\[
d\theta(t) = -\nabla L(\theta(t))\,dt + \Sigma(\theta(t))^{1/2}\,dW_t,
\]
where \(W_t\) is Brownian motion and \(\Sigma\) reflects gradient noise.

\medskip

\noindent
\textbf{Interpretation.}
Training is no longer a deterministic flow, but a noisy dynamical process. This noise can:
\begin{itemize}
    \item help escape unstable regions,
    \item bias trajectories toward some minima rather than others,
    \item influence both optimization and generalization.
\end{itemize}

\medskip

\noindent
Hence training dynamics are shaped by both geometry and stochasticity.

\section{Feature Learning vs Kernel-Like Regimes}

A central modern question is whether training mainly learns new features or behaves approximately like a linear model around initialization.

\subsection{Linearization Around Initialization}

Let \(\theta_0\) denote the initialization. For parameters \(\theta\) near \(\theta_0\), one may linearize the network:
\[
f_\theta(x)
\approx
f_{\theta_0}(x)
+
\nabla_\theta f_{\theta_0}(x)^\top (\theta-\theta_0).
\]

\medskip

\noindent
This approximation treats the model as locally linear in parameters. If training remains close to initialization, the behavior of the network may be well approximated by this linearized model.

\subsection{Neural Tangent Kernel}

The associated kernel is the \textbf{neural tangent kernel} (NTK):
\[
K(x,x')
=
\nabla_\theta f_{\theta_0}(x)^\top
\nabla_\theta f_{\theta_0}(x').
\]

\medskip

\noindent
\textbf{High-level meaning of the NTK.}
\begin{itemize}
    \item it converts the network into a kernel method at initialization,
    \item in suitable wide-network regimes, training becomes approximately linear,
    \item predictions evolve similarly to kernel regression with kernel \(K\).
\end{itemize}

\medskip

\noindent
In the infinite-width limit, this kernel often becomes deterministic, and the dynamics simplify considerably. This regime is sometimes called \textbf{lazy training}, because the parameters move little and the network behaves almost like its first-order Taylor expansion.

\subsection{Kernel-Like Regime}

In a kernel-like or lazy regime:
\begin{itemize}
    \item features do not change much during training,
    \item learning happens mainly through reweighting fixed random features,
    \item analysis is simpler and often mathematically tractable.
\end{itemize}

\medskip

\noindent
This regime is important because it connects deep learning to classical kernel methods and yields rigorous results in some asymptotic settings.

\subsection{Feature-Learning Regime}

In practical deep learning, however, networks often do more than linearized fitting.

\medskip

\noindent
\textbf{Feature learning} means that internal representations evolve significantly during training. Hidden layers reshape the geometry of the data so that later layers can solve the task more easily.

\medskip

\noindent
In this regime:
\begin{itemize}
    \item the tangent features themselves change,
    \item the network moves far enough from initialization that linearization becomes incomplete,
    \item representation learning becomes central to performance.
\end{itemize}

\medskip

\noindent
\textbf{Key contrast.}
\begin{itemize}
    \item \textbf{kernel-like view:} learning acts on nearly fixed features,
    \item \textbf{feature-learning view:} learning changes the features themselves.
\end{itemize}

\medskip

\noindent
A major theme of modern theory is understanding when each regime dominates and how they interact.

\section{Representation Geometry in Trained Networks}

Deep learning is not only about reducing loss. It also changes the internal geometry of data representations.

\subsection{Representation as Geometry}

For an intermediate layer \(h_\theta(x)\), inputs are mapped into a feature space. Training changes this map, and hence changes the geometry of the dataset in representation space.

\medskip

\noindent
\textbf{Representation geometry asks:}
\begin{itemize}
    \item which examples become closer or farther apart,
    \item which directions are expanded or compressed,
    \item how class structure becomes more separable,
    \item how invariances are formed.
\end{itemize}

\subsection{Separation and Flattening}

Empirically, trained networks often transform data so that:
\begin{itemize}
    \item examples from the same class cluster more tightly,
    \item examples from different classes become more separated,
    \item nuisance variations are suppressed,
    \item task-relevant directions are amplified.
\end{itemize}

\medskip

\noindent
This can be viewed as a geometric simplification of the learning problem.

\subsection{Feature Learning Geometry}

\textbf{Feature learning geometry} refers to how the network reshapes the data manifold during training.

\medskip

\noindent
At a high level, one may think of hidden layers as performing successive geometric operations:
\begin{itemize}
    \item disentangling factors of variation,
    \item flattening or straightening class manifolds,
    \item aligning representations with decision boundaries,
    \item creating coordinates in which simple classifiers work well.
\end{itemize}

\medskip

\noindent
This viewpoint helps explain why deep networks can outperform fixed-feature methods: they do not merely fit outputs; they actively reorganize the input geometry.

\subsection{Residual Networks and Dynamical Depth}

Residual networks provide an especially suggestive bridge between geometry and dynamics. A residual block has the form
\[
h_{k+1}=h_k + \Phi_k(h_k),
\]
which resembles a discretized differential equation.

\medskip

\noindent
Thus depth can be interpreted as time in a representation flow. Under this view:
\begin{itemize}
    \item layers evolve features gradually,
    \item deeper models implement longer geometric transformations,
    \item representation learning becomes a controlled dynamical process.
\end{itemize}

\medskip

\noindent
This is one reason dynamical-systems language is so natural in modern deep learning theory.

\section{A Unifying Perspective}

The geometry and dynamics of learning can be organized around several complementary viewpoints:

\begin{itemize}
    \item \textbf{loss landscapes:} the global and local structure of \(L(\theta)\),
    \item \textbf{curvature and Hessians:} second-order information shaping optimization,
    \item \textbf{gradient flow:} continuous-time evolution of parameters,
    \item \textbf{kernel-like regimes:} linearized training near initialization,
    \item \textbf{feature learning:} task-dependent evolution of representations,
    \item \textbf{representation geometry:} how the network restructures data internally.
\end{itemize}

\medskip

\noindent
\textbf{Core idea.}
Understanding deep learning requires studying both where optimization moves in parameter space and how those motions change the function and representation geometry seen by the data.

\medskip

\noindent
\textbf{Key insight.}
Modern learning theory is not only about approximation or statistics. It is also about geometry, dynamics, and the evolving structure of learned representations.

\section{Outlook}

This chapter developed a richer picture of learning by combining geometric and dynamical viewpoints:
\begin{itemize}
    \item parameter space versus function space,
    \item loss landscapes and high-dimensional critical structure,
    \item curvature, Hessian intuition, and optimization geometry,
    \item continuous-time gradient flow and stochastic dynamics,
    \item NTK and lazy-training perspectives,
    \item feature learning and representation geometry.
\end{itemize}

\medskip

\noindent
\textbf{Guiding principle.}
Learning is not merely the search for low loss. It is a structured dynamical process that moves through parameter space while reshaping the geometry of functions and representations.

\section*{Exercises}

\begin{enumerate}
    \item \textbf{Parameter space and function space.}
    Explain why two very different parameter vectors \(\theta\) and \(\theta'\) may correspond to similar functions \(f_\theta\) and \(f_{\theta'}\). Give at least two reasons this can happen in neural networks.

    \item \textbf{Quadratic approximation.}
    Starting from the second-order Taylor approximation
    \[
    L(\theta+v)\approx L(\theta)+\nabla L(\theta)^\top v+\frac12 v^\top H(\theta)v,
    \]
    explain the meaning of the Hessian term. What does it tell us about curvature in the direction \(v\)?

    \item \textbf{Eigenvalue interpretation.}
    Suppose the Hessian at a point has one negative eigenvalue, several large positive eigenvalues, and many near-zero eigenvalues. What does this suggest about the local geometry of the loss?

    \item \textbf{Simple gradient flow.}
    Consider the one-dimensional loss
    \[
    L(\theta)=\frac12 \theta^2.
    \]
    Write the gradient-flow equation and solve it explicitly. Show that the loss decreases monotonically along the trajectory.

    \item \textbf{Another gradient flow.}
    Consider
    \[
    L(\theta)=\frac14 \theta^4.
    \]
    Write the gradient-flow equation. Compare qualitatively its convergence near \(\theta=0\) with the quadratic case.

    \item \textbf{Fixed points and stability.}
    Let gradient flow satisfy
    \[
    \frac{d\theta}{dt}=-\nabla L(\theta).
    \]
    Why are critical points of \(L\) fixed points of the flow? How does curvature help distinguish stable equilibria from unstable ones?

    \item \textbf{Linearized training.}
    Derive the first-order Taylor approximation
    \[
    f_\theta(x)\approx f_{\theta_0}(x)+\nabla_\theta f_{\theta_0}(x)^\top(\theta-\theta_0).
    \]
    In what sense does this approximation lead to a kernel method?

    \item \textbf{NTK versus feature learning.}
    Compare the following two views:
    \begin{itemize}
        \item training as kernel-like regression with nearly fixed features,
        \item training as feature learning in which internal representations evolve.
    \end{itemize}
    What kinds of behavior can the second view explain that the first may miss?

    \item \textbf{Representation geometry.}
    Describe how training might change the geometry of hidden representations in a classification network. What does it mean for class manifolds to become more separable?

    \item \textbf{Residual networks and dynamics.}
    Consider the residual update
    \[
    h_{k+1}=h_k+\Phi_k(h_k).
    \]
    Explain why this resembles a discretized differential equation. What insight does this give into the role of depth?

    \item \textbf{Flat and sharp regions.}
    Why might flatness be related to robustness of a minimum? Why should one nevertheless be cautious about using flatness as a definitive explanation of generalization?

    \item \textbf{Synthesis.}
    Write a short essay comparing two pictures of deep learning:
    \begin{itemize}
        \item the optimization of a complicated loss landscape in parameter space,
        \item the gradual evolution of useful representations in function space.
    \end{itemize}
    Explain why both pictures are needed.
\end{enumerate}